\section{Découverte~$\boldsymbol{\alpha}$~: meilleur ajustement}
\label{sec:alpha}

\subsection{Contexte et source}
Ce candidat a été récupéré du \textbf{Cahier~1, chapitre~XV, page~19}. L'étape d'extraction par vision a classifié le contenu comme une \emph{fonction thêta mock} avec une confiance de 0,75.

\subsection{Eta-quotient découvert}
\begin{equation}
  f_\alpha(\tau) = q^{1/24} \cdot \eta(\tau)^{1} \, \eta(5\tau)^{-1} \, \eta(6\tau)^{2} \, \eta(7\tau)^{-1} \, \eta(8\tau)^{4} \, \eta(9\tau)^{-2} \, \eta(10\tau)^{12} \, \eta(11\tau)^{-3} \, \eta(12\tau)^{4}.
\end{equation}
\textbf{Énergie RAMA~:} $E = 1{,}1272$. \quad \textbf{Erreur d'ajustement~:} $I = 0{,}0705$ (7{,}05\,\%).

\subsection{Calcul exact des invariants}
\begin{proposition}[Invariants de $f_\alpha$]
\begin{align}
  \ceff &= \frac{3505}{1386} \approx 2{,}5289, \\
  k &= \frac{16}{2} = 8, \\
  p &= \frac{150}{24} = \frac{25}{4}.
\end{align}
\end{proposition}

\begin{proof}
Chaque calcul découle directement des définitions. Les sommes sont finies et exactes. La valeur rationnelle de $\ceff$ est confirmée par le noyau Lean~4.
\end{proof}

\subsection{Analyse de modularité}
\begin{proposition}
$\sum_d d \cdot r_d = 150 \not\equiv 0 \pmod{24}$. Par conséquent, $f_\alpha$ ne satisfait \textbf{pas} la condition~(M2) pour une forme modulaire holomorphe de niveau~1. C'est un eta-quotient \emph{généralisé} ou une composante d'une forme mock modulaire~\cite{zwegers2002,bringmann_ono_2006}.
\end{proposition}

\begin{proof}
$150 = 6 \times 24 + 6$, donc $150 \equiv 6 \pmod{24}$.
\end{proof}

\subsection{Développement en $q$-série}
Les premiers coefficients calculés par le script Python \texttt{compute\_alpha.py} sont~:
\[
  a(0..14) = [1, -1, -1, 0, 0, 2, -3, 3, -3, 5, -7, 7, 6, -9, 25].
\]
